\documentclass{article} % Clase base sin tamaño

% Página horizontal
\usepackage[a4paper,landscape,margin=2cm]{geometry}

% Arial y tamaño exacto
\usepackage{fontspec}
\setmainfont{Arial}
\usepackage{xcolor}
% Para tamaño grande explícito
\usepackage{setspace}
\usepackage{graphicx}

\usepackage{amsmath}
\usepackage[thmmarks, amsmath, thref]{ntheorem}
\usepackage{amsfonts}
\usepackage{amssymb}
\usepackage{polyglossia}       % Paquete multilingüe moderno
\setmainlanguage{spanish}   

\newcommand{\bigtable}[1]{{\fontsize{30pt}{36pt}\selectfont\renewcommand{\arraystretch}{2}#1}}


\begin{document}

% TODO lo que esté entre estas llaves tendrá tamaño 30pt real
{\fontsize{30pt}{36pt}\selectfont

\begin{center}
    Taller 1: Econometría 

Profesora:    Jocelyn Tapia Stefanoni\\
Ayudantes:   Matías Balboa y   Barbara Buffo
\end{center}



\textbf{Instrucciones:} 

\begin{itemize}
    \item Usted dispone de 60 minutos  para intentar los 100 puntos del taller.
    \item Cuide la presentación y redacción de sus respuestas.
    \item Puede utilizar su computador y los apuntes tanto de clase como de ayudantía.
    \item Descargue sus respuestas en un archivo con la extensión .ipynb. Debe enviar el enlace al notebook en Google Colab junto con el archivo .ipynb adjunto por correo electrónico a la profesora (jocelyn.tapias@usm.cl), con copia a la ayudante (econometriausm482@gmail.com). 
    \end{itemize}

    Para este taller, utilice los datos del archivo \texttt{wage2} del paquete wooldridge. 


\begin{enumerate}
\item (10 puntos) Describa la base de datos, identificando su origen y las variables que contienen. 
    \item (10 puntos) Calcule y presente estadísticos descriptivos relevantes para todas las variables contenidas en la base de datos.

\item (10 puntos) Obtenga e interprete la matriz de correlaciones entre las variables incluidas en la base de datos.

\item (10 puntos) Estime una regresión lineal simple del \texttt{IQ} sobre \texttt{educ}, obteniendo el coeficiente de la pendiente, que denotará como $\tilde{\delta}_1$.

\item (10 puntos) Estime una regresión lineal simple del logaritmo del salario $\log(\texttt{wage})$ sobre \texttt{educ}, obteniendo el coeficiente de la pendiente, que denotará como $\tilde{\beta}_1$.

\item (10 puntos) Estime una regresión lineal múltiple del logaritmo del salario $\log(\texttt{wage})$ sobre \texttt{educ} y \texttt{IQ}, obteniendo los coeficientes de las pendientes, que denotará como $\hat{\beta}_1$ y $\hat{\beta}_2$, respectivamente.

\item (14 puntos) Verifique empíricamente la siguiente relación:
\[
\tilde{\beta}_1 \approx \hat{\beta}_1 + \hat{\beta}_2 \cdot \tilde{\delta}_1
\]
donde $\tilde{\delta}_1$ es obtenido del punto (4), $\tilde{\beta}_1$ del punto (5), y $\hat{\beta}_1$ y $\hat{\beta}_2$ del punto (6).

\item (16 puntos) Determine e interprete la variabilidad explicada y la variabilidad no explicada 
 por el modelo múltiple estimado en el punto (6).

 \item (10 puntos) Determine el error estándar de los estimadores de MCO del modelo del punto (6). 
%\item (22 puntos) Sugiera y justifique claramente qué variables adicionales podría agregar al modelo para mejorar su especificación. Incorpórelas, estime nuevamente el modelo, y compare e interprete sus resultados en relación al modelo original estimado en el punto (6).
\end{enumerate}

\end{document}